\section{Introduction}
Systematic reviews are essential tools in evidence-based medicine, providing comprehensive and unbiased summaries of scientific knowledge across medicine and social sciences. At the heart of these reviews lies a deceptively complex task: the formulation of Boolean search queries capable of retrieving all relevant literature (high recall) without overburdening researchers with non-relevant results (high precision). A well-formulated Boolean query directly affects the cost, efficiency, and reproducibility of the review process.

Large language models (LLMs) have been explored as tools for automatically generating Boolean queries from an initial research question or topic~\cite{wang2023can, wang2025reassessing,staudinger2024reproducibility}. While conceptually easy to use, prompt-based LLM methods have shown major limitations---often retrieving far fewer relevant studies than expert-crafted queries (low recall), well below the thresholds typically required for systematic reviews (e.g., 10--40\% recall instead of 80--90\%)~\cite{wang2025reassessing}. These limitations highlight the need for new methods that go beyond zero-shot prompting and instead are optimized to generate queries based on retrieval effectiveness.

A natural alternative to prompting is supervised fine-tuning on example queries. However, this is impractical for Boolean query generation, where no single high-quality ground-truth query exists. Expert-crafted queries are often inconsistent, and sub-optimal~\cite{scells2018generating}. Further, existing datasets are too small: fewer than 200 training pairs, many sourced from Cochrane~\cite{kanoulas2018clef, wang2022seed}.%
\footnote{Cochrane does not support direct API access; prior work translates these queries into MEDLINE format for PubMed execution~\cite{wang2025reassessing}.}

We propose a reinforcement learning (RL) framework called \textbf{AutoBool} to train LLMs for Boolean query generation, enabling direct optimization for retrieval effectiveness. Rather than relying on handcrafted prompts or static templates, our model learns to balance recall and precision through feedback from document retrieval. To support this task, we construct a large-scale training dataset of \num{32794} systematic reviews mined from the PubMed Central (PMC) Open Access corpus, and introduce a retrieval-grounded reward function aligned with the screening goals of systematic review creation.
We further release a large-scale evaluation dataset comprising \num{32794} topics with high-quality relevance labels---an order of magnitude larger than prior benchmarks such as CLEF TAR and the Seed Collection,%
\footnote{Which contain 118 and 40 topics overall respectively.} 
and with a smaller PubTemp set designed to be free from data leakage. Figure~\ref{fig:overview} illustrates our complete pipeline, from dataset construction to RL training.

Our experiments show that \textbf{AutoBool}-trained models substantially outperform prompt-based zero-shot baselines, narrow the gap with expert-crafted queries, and even exceed the effectiveness of much larger commercial LLMs such as GPT-4o and O3 in high-recall retrieval scenarios, despite relying on significantly smaller backbones.
